% GENERATED FILE. Edit canonical lessons, then run npm run generate:books.
% canonical-source-hash: fnv1a64:6ced614e6a2ff51e
% canonical-lessons: ML-C506-suryodayam, ML-C506-suryastamayam, ML-C506-cakravalam, ML-C506-antariksam, ML-C506-malakkar

\chapter{The Sky}
\label{ch:ml-the-sky}
\hlchaptermodality{506}

\begin{chapteropening}
\textbf{By the end of this chapter:} \emph{I can say sunrise, sunset, the horizon, the atmosphere and rain clouds.}

Complete the last lesson of chapter 506: I can say sunrise, sunset, the horizon, the atmosphere and rain clouds.
\end{chapteropening}

\section[\texorpdfstring{\ml{സൂര്യോദയം} (sūryōdayaṁ)}{സൂര്യോദയം (sūryōdayaṁ)}]{\ml{സൂര്യോദയം} (sūryōdayaṁ) — sunrise}

\label{lesson:ML-C506-suryodayam}



\begin{quote}
Before the new one: say the Malayalam for global warming, then the Malayalam for a rain pit.
\end{quote}

\subsection*{You'll want to know: \ml{സൂര്യോദയം}}
\textbf{\ml{സൂര്യോദയം}} — \emph{sūryōdayaṁ} — \textquotedblleft{}sunrise\textquotedblright{}.

\begin{quote}
\ml{സൂര്യോദയം} \ml{കാണാൻ} \ml{ഞങ്ങൾ} \ml{പുലർച്ചെ} \ml{എഴുന്നേറ്റു}.
\end{quote}

\begin{quote}
\emph{sūryōdayaṁ kāṇān ñaṅṅaḷ pularcce eḻunnēṟṟu.}
\end{quote}

\begin{quote}
We got up at dawn to watch the sunrise.
\end{quote}

\ml{സൂര്യൻ}, the sun, plus \ml{ഉദയം}, rising: the vowels join as -\ml{ോ}-.

\begin{cousinweb}
\textbf{\ml{സൂര്യോദയം}} (\emph{sūryōdayaṁ}) is borrowed from Sanskrit \emph{sūryodaya}, sunrise, from \emph{sūrya}, \textquotedblleft{}the sun\textquotedblright{}, and \emph{udaya}, \textquotedblleft{}rising\textquotedblright{} (Monier-Williams).
\end{cousinweb}

\begin{grammarlens}[title={What you've built}]
One more word for getting around.
\end{grammarlens}

\subsection*{Your turn}
\begin{itemize}
\raggedright
  \item \emph{Say it:} \emph{sūryōdayaṁ}
  \item \emph{Say it:} \emph{sūryōdayaṁ}, once more
  \item \emph{Say it:} \emph{sūryōdayaṁ kāṇān ñaṅṅaḷ pularcce eḻunnēṟṟu.} — We got up at dawn to watch the sunrise.
  \item \emph{Recall:} say the Malayalam for global warming, then the Malayalam for a rain pit, then say \emph{sūryōdayaṁ} again
\end{itemize}

\begin{tcolorbox}[breakable,colback=teal!4,colframe=teal!35!black,title={Before you move on}]
What does \ml{സൂര്യോദയം} mean? (\textquotedblleft{}Sunrise\textquotedblright{}.) Say it once more.
\end{tcolorbox}
\section[\texorpdfstring{\ml{സൂര്യാസ്തമയം} (sūryāstamayaṁ)}{സൂര്യാസ്തമയം (sūryāstamayaṁ)}]{\ml{സൂര്യാസ്തമയം} (sūryāstamayaṁ) — sunset}

\label{lesson:ML-C506-suryastamayam}



\begin{quote}
Before the new one: say the Malayalam for a rain pit, then the Malayalam for sunrise.
\end{quote}

\subsection*{You'll want to know: \ml{സൂര്യാസ്തമയം}}
\textbf{\ml{സൂര്യാസ്തമയം}} — \emph{sūryāstamayaṁ} — \textquotedblleft{}sunset\textquotedblright{}.

\begin{quote}
\ml{കടപ്പുറത്ത്} \ml{സൂര്യാസ്തമയം} \ml{വളരെ} \ml{സുന്ദരമായിരുന്നു}.
\end{quote}

\begin{quote}
\emph{kaṭappuṟattŭ sūryāstamayaṁ vaḷare sundaramāyirunnu.}
\end{quote}

\begin{quote}
The sunset on the beach was very beautiful.
\end{quote}

\ml{സൂര്യൻ} plus \ml{അസ്തമയം}, setting. You met the verb \ml{അസ്തമിക്കുക}, to set.

\begin{cousinweb}
\textbf{\ml{സൂര്യാസ്തമയം}} (\emph{sūryāstamayaṁ}) is a Malayalam compound built from Sanskrit \emph{sūrya}, \textquotedblleft{}the sun\textquotedblright{}, and \emph{astamaya}, \textquotedblleft{}setting of the sun\textquotedblright{} (Monier-Williams).
\end{cousinweb}

\begin{grammarlens}[title={What you've built}]
One more word for getting around.
\end{grammarlens}

\subsection*{Your turn}
\begin{itemize}
\raggedright
  \item \emph{Say it:} \emph{sūryāstamayaṁ}
  \item \emph{Say it:} \emph{sūryāstamayaṁ}, once more
  \item \emph{Say it:} \emph{kaṭappuṟattŭ sūryāstamayaṁ vaḷare sundaramāyirunnu.} — The sunset on the beach was very beautiful.
  \item \emph{Recall:} say the Malayalam for a rain pit, then the Malayalam for sunrise, then say \emph{sūryāstamayaṁ} again
\end{itemize}

\begin{tcolorbox}[breakable,colback=teal!4,colframe=teal!35!black,title={Before you move on}]
What does \ml{സൂര്യാസ്തമയം} mean? (\textquotedblleft{}Sunset\textquotedblright{}.) Say it once more.
\end{tcolorbox}
\section[\texorpdfstring{\ml{ചക്രവാളം} (cakravāḷaṁ)}{ചക്രവാളം (cakravāḷaṁ)}]{\ml{ചക്രവാളം} (cakravāḷaṁ) — the horizon}

\label{lesson:ML-C506-cakravalam}



\begin{quote}
Before the new one: say the Malayalam for sunrise, then the Malayalam for sunset.
\end{quote}

\subsection*{You'll want to know: \ml{ചക്രവാളം}}
\textbf{\ml{ചക്രവാളം}} — \emph{cakravāḷaṁ} — \textquotedblleft{}the horizon\textquotedblright{}.

\begin{quote}
\ml{ചക്രവാളത്തിൽ} \ml{ഒരു} \ml{കപ്പൽ} \ml{കണ്ടു}.
\end{quote}

\begin{quote}
\emph{cakravāḷattil oru kappal kaṇṭu.}
\end{quote}

\begin{quote}
A ship was seen on the horizon.
\end{quote}

The circle where the sky seems to meet the land or the sea.

\begin{cousinweb}
\textbf{\ml{ചക്രവാളം}} (\emph{cakravāḷaṁ}) is borrowed from Sanskrit \emph{cakravāla}, \textquotedblleft{}a circle\textquotedblright{} (Monier-Williams); in Malayalam it names the horizon.
\end{cousinweb}

\begin{grammarlens}[title={What you've built}]
One more word for getting around.
\end{grammarlens}

\subsection*{Your turn}
\begin{itemize}
\raggedright
  \item \emph{Say it:} \emph{cakravāḷaṁ}
  \item \emph{Say it:} \emph{cakravāḷaṁ}, once more
  \item \emph{Say it:} \emph{cakravāḷattil oru kappal kaṇṭu.} — A ship was seen on the horizon.
  \item \emph{Recall:} say the Malayalam for sunrise, then the Malayalam for sunset, then say \emph{cakravāḷaṁ} again
\end{itemize}

\begin{tcolorbox}[breakable,colback=teal!4,colframe=teal!35!black,title={Before you move on}]
What does \ml{ചക്രവാളം} mean? (\textquotedblleft{}The horizon\textquotedblright{}.) Say it once more.
\end{tcolorbox}
\section[\texorpdfstring{\ml{അന്തരീക്ഷം} (antarīkṣaṁ)}{അന്തരീക്ഷം (antarīkṣaṁ)}]{\ml{അന്തരീക്ഷം} (antarīkṣaṁ) — the atmosphere}

\label{lesson:ML-C506-antariksam}



\begin{quote}
Before the new one: say the Malayalam for sunset, then the Malayalam for the horizon.
\end{quote}

\subsection*{You'll want to know: \ml{അന്തരീക്ഷം}}
\textbf{\ml{അന്തരീക്ഷം}} — \emph{antarīkṣaṁ} — \textquotedblleft{}the atmosphere\textquotedblright{}.

\begin{quote}
\ml{മഴയ്ക്ക്} \ml{ശേഷം} \ml{അന്തരീക്ഷം} \ml{തണുത്തു}.
\end{quote}

\begin{quote}
\emph{maḻaykkŭ śēṣaṁ antarīkṣaṁ taṇuttu.}
\end{quote}

\begin{quote}
After the rain the air grew cool.
\end{quote}

Also the mood of a place: a calm atmosphere is \ml{ശാന്തമായ} \ml{അന്തരീക്ഷം}.

\begin{cousinweb}
\textbf{\ml{അന്തരീക്ഷം}} (\emph{antarīkṣaṁ}) is borrowed from Sanskrit \emph{antarikṣa}, \textquotedblleft{}the space between heaven and earth\textquotedblright{} (Monier-Williams).
\end{cousinweb}

\begin{grammarlens}[title={What you've built}]
One more word for getting around.
\end{grammarlens}

\subsection*{Your turn}
\begin{itemize}
\raggedright
  \item \emph{Say it:} \emph{antarīkṣaṁ}
  \item \emph{Say it:} \emph{antarīkṣaṁ}, once more
  \item \emph{Say it:} \emph{maḻaykkŭ śēṣaṁ antarīkṣaṁ taṇuttu.} — After the rain the air grew cool.
  \item \emph{Recall:} say the Malayalam for sunset, then the Malayalam for the horizon, then say \emph{antarīkṣaṁ} again
\end{itemize}

\begin{tcolorbox}[breakable,colback=teal!4,colframe=teal!35!black,title={Before you move on}]
What does \ml{അന്തരീക്ഷം} mean? (\textquotedblleft{}The atmosphere\textquotedblright{}.) Say it once more.
\end{tcolorbox}
\section[\texorpdfstring{\ml{മഴക്കാർ} (maḻakkār)}{മഴക്കാർ (maḻakkār)}]{\ml{മഴക്കാർ} (maḻakkār) — rain clouds}

\label{lesson:ML-C506-malakkar}



\begin{quote}
Before the new one: say the Malayalam for the horizon, then the Malayalam for the atmosphere.
\end{quote}

\subsection*{You'll want to know: \ml{മഴക്കാർ}}
\textbf{\ml{മഴക്കാർ}} — \emph{maḻakkār} — \textquotedblleft{}rain clouds\textquotedblright{}.

\begin{quote}
\ml{ആകാശത്ത്} \ml{മഴക്കാർ} \ml{നിറഞ്ഞു}.
\end{quote}

\begin{quote}
\emph{ākāśattŭ maḻakkār niṟaññu.}
\end{quote}

\begin{quote}
The sky filled with rain clouds.
\end{quote}

\ml{കാർ} is a dark cloud: \ml{മഴ}, rain, plus \ml{കാർ}.

\begin{cousinweb}
\textbf{\ml{മഴക്കാർ}} (\emph{maḻakkār}) is a native Dravidian compound of \ml{മഴ}, rain (DEDR 4753), and \ml{കാർ}, darkness or a dark cloud (DEDR 1278).
\end{cousinweb}

\begin{grammarlens}[title={What you've built}]
One more word for getting around.
\end{grammarlens}

\subsection*{Your turn}
\begin{itemize}
\raggedright
  \item \emph{Say it:} \emph{maḻakkār}
  \item \emph{Say it:} \emph{maḻakkār}, once more
  \item \emph{Say it:} \emph{ākāśattŭ maḻakkār niṟaññu.} — The sky filled with rain clouds.
  \item \emph{Recall:} say the Malayalam for the horizon, then the Malayalam for the atmosphere, then say \emph{maḻakkār} again
\end{itemize}

\begin{tcolorbox}[breakable,colback=teal!4,colframe=teal!35!black,title={Before you move on}]
What does \ml{മഴക്കാർ} mean? (\textquotedblleft{}Rain clouds\textquotedblright{}.) Say it once more.
\end{tcolorbox}
